\begin{titlepage}
\color{navy}\sffamily\raggedright\vspace*{30pt}
{\large THE ENERGY BALANCE PROJECT\par}\vspace{30pt}
{\fontsize{32}{38}\selectfont\bfseries Implementation, Preservation and Evidence\par}
\vspace{20pt}{\Large Part III\par}\vspace{12pt}
{\large Revised Gaussian Explanatory Edition\par}\vfill
{\large Konrad Grzyb\par}\vspace{12pt}
{\normalsize Author-review manuscript\par September 2026\par}
\end{titlepage}
\chapter*{What exists, and what is proposed}
\addcontentsline{toc}{chapter}{What exists, and what is proposed}
This volume explains how a mathematical account becomes a testable scientific
system. It also explains how to avoid reporting a proposed system as though
it already existed. Those are inseparable tasks: the credibility of an
implementation depends on knowing exactly what it implements and what its
tests actually establish.

The project has historical implementations and experiments under earlier
potential and controller definitions. It also has a prospective Local
Gaussian programme, preserved candidate code, and source-locked framework
components on other Git lineages. These are not one interchangeable engine.
The active chapters identify their differences. A historical appendix retains
the original D1--D10, quote-validation, adversarial and service-alignment
chapters without rewriting their results under the Gaussian formula.

\begin{lawbox}{The central reading rule}
In this book, a prospective interface is not an implemented service; a passing
arithmetic fixture is not a scientific campaign; and a historical result is
evidence only for its own registered model and conditions.
\end{lawbox}

This edition follows Part II's revised numbering, beginning at Chapter 47.
Original chapter numbers in the historical appendix remain archival labels.
The supplied original Part III PDF is preserved. New chapters are editable
reconstructed source, retaining the original teaching style rather than
claiming access to the original typesetting manuscript.

No Gaussian runtime, scientific trajectory or new experiment was executed to
write this book. The prospective chapters describe responsibilities and
acceptance boundaries, not accepted runnable APIs or a sealed study packet.
Parameters, thresholds, seed families, horizons and cloud configurations are
not selected here. Historical results are not erased because their model is
no longer the active candidate.

\chapter*{Reading the software and the evidence}
\addcontentsline{toc}{chapter}{Reading the software and the evidence}
Chapters 47--55 follow one proposed scientific event from its declared state
through forcing, candidates, physical screening, valuation, affordability,
selection, execution records and reproducibility. The ordering is the new
programme's intended profile; reconciliation with the existing framework
remains an explicit dependency, not a completed migration.

Chapters 56--57 explain the historical architecture and evidence. They are
the bridge to the unaltered historical pages at the back. Chapter 58 examines
measurement and the open loss-accounting boundary. Chapter 59 asks what a
long-run result could mean. Chapter 60 supplies a worked audit and a precise
map of the next separately authorized stages.

A reader interested in implementation should ask three questions of every
module: which facts may it read, which facts may it write, and what exact
claim does its output support? A reader interested in evidence should add:
which version ran, what comparison was fixed beforehand, and where are the
complete records, including refusals and negative outcomes?

The historical appendix reproduces publication-era prose. Words such as
``current'' there refer to that source's time, not to this revision. The later
finalized Gate 1D-C record is treated separately in Chapter 57. Neither the
appendix nor that later record demonstrates success of the new Gaussian
random-affordable programme.
